% Часть 3: димер и чистая дефазировка (PRB).

\subsection{Димер (PRB, ур.~24--34, рис.~2)}
\label{sec:kd_dimer}

\paragraph{Зачем.} На рис.~2(b) статьи кривые для 2D и 3D вдвое ниже, чем дают
формулы самой статьи (ур.~31--32) и наш код. Выведем формулы и проверим их
физическим пределом.

\paragraph{Система.} Два спина $1/2$ на расстоянии $d$, $K\equiv K_d(k_Fd)$. Синглет и
триплет:
\[
|1,1\rangle=|\!\uparrow\uparrow\rangle,\quad
|1,0\rangle=\tfrac{1}{\sqrt2}(|\!\uparrow\downarrow\rangle+|\!\downarrow\uparrow\rangle),\quad
|1,-1\rangle=|\!\downarrow\downarrow\rangle,\quad
|0,0\rangle=\tfrac{1}{\sqrt2}(|\!\uparrow\downarrow\rangle-|\!\downarrow\uparrow\rangle).
\]
Матрица ядра: $K_{11}=K_{22}=1$, $K_{12}=K_{21}=K$. Для вектора матричных элементов
$x=(x_1,x_2)$ на первом и втором спине
\[
\sum_{ll'}K_{ll'}x_lx_{l'}^*=|x_1|^2+|x_2|^2+2K\,\mathrm{Re}(x_1x_2^*).
\]

\paragraph{Переход $|1,0\rangle\to|0,0\rangle$ (рис.~2b).} Действуем $S^z_1$:
\[
S^z_1|1,0\rangle=\tfrac12\cdot\tfrac1{\sqrt2}(|\!\uparrow\downarrow\rangle-|\!\downarrow\uparrow\rangle)=\tfrac12|0,0\rangle,\qquad
S^z_2|1,0\rangle=-\tfrac12|0,0\rangle .
\]
$S^x$, $S^y$ меняют полный $S^z$ и на $|0,0\rangle$ элементов не имеют. Итак
$x=(\tfrac12,-\tfrac12)$ только для $a=z$:
\[
\sum_{ll'}K_{ll'}x_lx_{l'}=\tfrac14+\tfrac14+2K\cdot\big(-\tfrac14\big)=\tfrac12(1-K).
\]
Делим на $\tfrac12$ --- ту же сумму для одиночного спина (раздел~\ref{sec:kd_rate}).

\paragraph{Переход $|1,0\rangle\to|1,-1\rangle$ в сильном поле (рис.~2c).} $S^-_l$ опускает
спин $l$: $S^-_{1,2}|1,0\rangle=\tfrac1{\sqrt2}|\!\downarrow\downarrow\rangle$. С $S^x=\tfrac12(S^++S^-)$ и
$S^y=\tfrac1{2i}(S^+-S^-)$ для обоих $l$
\[
\langle1,-1|S^x_l|1,0\rangle=\tfrac1{2\sqrt2},\qquad \langle1,-1|S^y_l|1,0\rangle=\tfrac{i}{2\sqrt2}.
\]
Для $a=x$: $\tfrac18+\tfrac18+2K\cdot\tfrac18=\tfrac14(1+K)$; для $a=y$ то же; всего $\tfrac12(1+K)$.
\begin{resbox}[Скорости димера]
\[
\Gamma_{|1,0\rangle\to|0,0\rangle}=\Gamma_0\,(1-K_d(k_Fd)),\qquad
\Gamma_{|1,0\rangle\to|1,-1\rangle}=\Gamma_0\,(1+K_d(k_Fd)).
\]
\end{resbox}
\begin{paperbox}{PRB, ур.~28, 30 (1D) и ур.~31--34 (2D, 3D)}
\[
\Gamma_{|1,S_Z\rangle\to|0,0\rangle}=\Gamma_0(\Delta_{ST})\big[1-\cos^2(k_Fd)\big],\qquad
\Gamma_{|1,-1\rangle\to|1,0\rangle}=\Gamma_0(g\mu_BB_z)\big[1+\cos^2(k_Fd)\big],
\]
\[
\Gamma^d_{|1,S_Z\rangle\to|0,0\rangle}=\Gamma_0(\Delta_{S,T})\,F^d(k_Fd),\qquad
F^d(x)=\begin{cases}1-J_0^2(x),&d=2\\1-\mathrm{sinc}^2(x),&d=3\end{cases}
\]
\[
\Gamma^d_{|1,-1\rangle\to|1,0\rangle}=\Gamma_0(g\mu_BB_z)\,F^d_T(k_Fd),\qquad
F^d_T(x)=\begin{cases}1+J_0^2(x),&d=2\\1+\mathrm{sinc}^2(x),&d=3\end{cases}
\]
\end{paperbox}
\begin{cmpbox}
Формулы статьи совпадают с нашими во всех размерностях. Расхождение только
с напечатанным рисунком~2(b) --- ниже.
\end{cmpbox}

\paragraph{Физический смысл.} При $k_Fd=0$ ($K=1$) переход в синглет запрещён
($1-K=0$), а переход внутри триплета вдвое быстрее одиночного спина ($1+K=2$).
Это аналог суб- и сверхизлучения: амплитуды рассеяния на двух атомах
складываются в противофазе или в фазе.

\paragraph{Предел далёких спинов --- решающая проверка.} В 2D и 3D $K\to0$ при
$k_Fd\to\infty$. Далёкие спины рассеивают электроны независимо, и скорость
равна некогерентной сумме $\sum_l\sum_a|\langle f|S^a_l|i\rangle|^2$. Для обоих
переходов она равна $\tfrac12$ ($\tfrac14+\tfrac14$ и $\tfrac18\cdot4$) --- как для одиночного спина
(проверено численно).
\begin{resbox}[Обязательный предел]
\[
\Gamma_{|1,0\rangle\to|0,0\rangle},\ \Gamma_{|1,0\rangle\to|1,-1\rangle}\ \xrightarrow[k_Fd\to\infty]{}\ \Gamma_0\qquad(d=2,3).
\]
\end{resbox}
\begin{paperbox}{PRB, текст после ур.~34}
В пределе, где осцилляции затухли, субрадиантное (синглетное) состояние
излучает со скоростью, вдвое меньшей, чем одиночный спин.
\end{paperbox}
\begin{cmpbox}[Сравнение с рис.~2]
\begin{itemize}
\item Рис.~2(c): кривые 2D и 3D стремятся к 1, как положено.
\item Рис.~2(b): кривые 2D и 3D стремятся к 0.5, хотя обе панели нормированы
  на один и тот же $\Gamma_0$.
\item Оцифровка из PDF: кривые 2D и 3D рис.~2(b) совпадают с $(1-K)/2$ с
  точностью 0.001--0.003; кривая 1D рис.~2(b) и все кривые рис.~2(c) --- с
  формулами без множителя $\tfrac12$.
\end{itemize}
\end{cmpbox}

\paragraph{Откуда могла взяться «половина».} В 1D $K=\cos^2(k_Fd)$ не затухает, а
осциллирует со средним $\tfrac12$, поэтому среднее значение 1D-кривой $1-K$ равно
$\tfrac12$. Вероятно, этот вывод для 1D был перенесён на 2D и 3D, где он неверен.

\begin{cmpbox}[Вывод по рис.~2]
Ур.~31--32, наш расчёт, панель (c) и предел независимых спинов согласованы
между собой; кривые 2D и 3D панели (b) отличаются от них множителем $\tfrac12$.
Расхождение --- в рисунке статьи.
\end{cmpbox}

\subsection{Чистая дефазировка (PRB, ур.~15--16, 38--39, рис.~4, приложение C)}
\label{sec:kd_deph}

\paragraph{Зачем.} Рис.~4 статьи (дефазировка изинговских цепочек) не совпадает с
нашим расчётом. Оцифровка показала, что рисунок построен по ур.~16 в
напечатанном виде. Выведем дефазировку и проверим ур.~16.

\paragraph{Откуда берётся дефазировка.} Часть оператора коллапса с $\omega=0$
(раздел~\ref{sec:kd_davies}) диагональна:
\[
L_0=\sqrt{c}\sum_Ma_M|M\rangle\langle M|,\qquad a_M=\langle M|A|M\rangle,\qquad
c=\tfrac{\pi(\rho J)^2}{4}\lambda\,G(0)=\tfrac{\pi(\rho J)^2}{4}\lambda\,k_BT .
\]
Действие слагаемых Линдблада на когерентность $\rho_{MM'}=\langle M|\rho|M'\rangle$:
\[
\langle M|L_0\rho L_0^\dagger|M'\rangle=c\,a_Ma_{M'}\rho_{MM'},\qquad
\langle M|\tfrac12\{L_0^\dagger L_0,\rho\}|M'\rangle=\tfrac c2\big(a_M^2+a_{M'}^2\big)\rho_{MM'} .
\]
Разность (для вещественных $a$):
\[
c\Big(a_Ma_{M'}-\tfrac12a_M^2-\tfrac12a_{M'}^2\Big)\rho_{MM'}=-\tfrac c2(a_M-a_{M'})^2\rho_{MM'} .
\]
Когерентность затухает со скоростью $\tfrac c2(a_M-a_{M'})^2$, населённости не
меняются. Множитель $\tfrac12$ --- из структуры лиувиллиана. Суммируя по модам ядра:
\begin{resbox}[Скорость чистой дефазировки]
\begin{equation}
\gamma^{\rm ad}_{MM'}=\frac{\pi(\rho J)^2}{8}\,k_BT\sum_a\sum_{nn'}
K_d(k_F|\mathbf r_n-\mathbf r_{n'}|)\,\Delta S^a_n\,\Delta S^a_{n'},
\qquad \Delta S^a_n=S^a_{MM}(n)-S^a_{M'M'}(n).
\label{eq:kd_ad}
\end{equation}
До усреднения по поверхности Ферми здесь стоит
$\big|\sum_ne^{i\mathbf q\cdot\mathbf r_n}\Delta S_n\big|^2$: у $S_{MM}(n)$ и $S_{M'M'}(n)$ \emph{один и тот же}
фазовый множитель $e^{i\mathbf q\cdot\mathbf r_n}$, так как это элементы одного
оператора~\eqref{eq:kd_V}.
\end{resbox}
\begin{paperbox}{PRB, ур.~38--39}
\[
\gamma^{\rm ad}_{M,M'}=\frac{\pi(\rho J)^2}{8\hbar}\,k_BT\,\chi^{\rm ad}_{MM'},\qquad
\chi^{\rm ad}_{MM'}=\sum_a\big|S^a_{MM}-S^a_{M'M'}\big|^2;
\]
\[
\gamma^{\rm ad}_{M,M'}=\frac{\pi(\rho J)^2}{2\hbar}\,k_BT\,S^2\qquad (\text{состояния }\pm S).
\]
\end{paperbox}
\begin{cmpbox}
Ур.~38 --- это~\eqref{eq:kd_ad} для одного спина. Для состояний $\pm S$:
$\Delta S^z=2S$, и $\tfrac{\pi(\rho J)^2}{8}k_BT\cdot4S^2=\tfrac{\pi(\rho J)^2}{2}k_BTS^2$ --- ур.~39. Совпадает.
\end{cmpbox}

\paragraph{Три проверки формулы~\eqref{eq:kd_ad}.}
\begin{enumerate}
\item \textbf{Не зависит от начала координат.} Сдвиг всех атомов на $\mathbf R$
  умножает $\sum_ne^{i\mathbf q\cdot\mathbf r_n}\Delta S_n$ на общую фазу $e^{i\mathbf q\cdot\mathbf R}$, модуль не
  меняется.
\item \textbf{Один спин} даёт ур.~39 (выше).
\item \textbf{Изинговская цепочка.} Два основных состояния с $S^z_n=\pm s_nS$,
  $s_n=\pm1$: $\Delta S_n=2s_nS$, и
  \[
  \gamma^{\rm ad}=\tfrac{\pi(\rho J)^2}{8}k_BTS^2\cdot\Lambda,\qquad \Lambda=4\,s^TKs .
  \]
  Это и считает код.
\end{enumerate}

\paragraph{Что напечатано в ур.~16.}
\begin{paperbox}{PRB, ур.~15--16}
\[
\gamma^{\rm ad}_{M,M'}=\frac{\pi J^2}{2\hbar}\sum_{\mathbf k\mathbf k'}f(\epsilon_{\mathbf k})\big[1-f(\epsilon_{\mathbf k'})\big]\,
\chi^{\rm ad}_{M,M'}(\mathbf k-\mathbf k')\,\delta(\epsilon_{\mathbf k}-\epsilon_{\mathbf k'}),
\]
\[
\chi^{\rm ad}_{M,M'}(\mathbf q)=\sum_a\Big|\sum_{n=1}^{N}e^{i\mathbf q\cdot\mathbf r_n}S^a_{MM}(n)-e^{-i\mathbf q\cdot\mathbf r_n}S^a_{M'M'}(n)\Big|^2 .
\]
\end{paperbox}
Второе слагаемое под модулем записано с $e^{-i\mathbf q\cdot\mathbf r_n}$ вместо $e^{+i\mathbf q\cdot\mathbf r_n}$. К
чему это приводит:
\begin{itemize}
\item Один спин в точке $\mathbf r$, состояния $\pm S$:
  \[
  |e^{i\mathbf q\cdot\mathbf r}S-e^{-i\mathbf q\cdot\mathbf r}(-S)|^2=4S^2\cos^2(\mathbf q\cdot\mathbf r)
  \ \xrightarrow{\text{среднее}}\ 2S^2\big[1+K_d(2k_Fr)\big]
  \]
  (через $\cos^2y=\tfrac12(1+\cos2y)$). Дефазировка одиночного спина зависела бы от
  его положения относительно произвольного начала координат и совпадала
  бы с ур.~39 только при $r=0$.
\item Цепочка: та же замена даёт слагаемое с суммой координат,
  \[
  \Lambda=2\sum_{ll'}s_ls_{l'}\big[K(k_F|r_l-r_{l'}|)+K(k_F|r_l+r_{l'}|)\big].
  \]
\end{itemize}
\begin{paperbox}{PRB, приложение C: ур.~C1, C3, C4}
\[
\chi^{\rm ad}_{MM'}(\mathbf q)=4S^2\sum_{ll'}\cos(\mathbf q\cdot\mathbf r_l)\cos(\mathbf q\cdot\mathbf r_{l'}),
\]
\[
\Lambda^{FM}_{1D}(\xi,N)=2\big[N^2+\csc^2\xi\,\sin^2(\xi N)\cos^2(\xi(N+1))\big],
\]
\[
\Lambda^{FM}_{3D}(\xi,N)=2\sum_{ll'}\big[\mathrm{sinc}^2(\xi|l-l'|)+\mathrm{sinc}^2(\xi|l-l'|)\big],\qquad \xi=k_Fd .
\]
\end{paperbox}
\begin{cmpbox}[Сравнение с приложением C и рис.~4]
\begin{itemize}
\item Ур.~C1 --- это ур.~16 с знаком $e^{-i\mathbf q\cdot\mathbf r}$ (произведение косинусов), то есть
  форма, зависящая от начала координат.
\item По этой форме построены ур.~C1, C3, C5, C6 и рис.~4. Оцифровка
  рисунка: все девять кривых совпадают с ней (медиана отклонения
  0.005--0.018 на шкале 0--4) и не совпадают с~\eqref{eq:kd_ad} (0.08--0.50).
\item Ур.~C4 напечатано с $|l-l'|$ в обоих слагаемых, то есть совпадает
  с~\eqref{eq:kd_ad}, а не с рисунком.
\item Ур.~C6 (нечётное $N$) не совпадает даже с суммой ур.~C1. Сумма по C1:
  \[
  \tfrac12\Big[4+\Big(1+\frac{\cos((2N+1)\xi)}{\cos\xi}\Big)^2\Big].
  \]
\end{itemize}
\end{cmpbox}

\begin{cmpbox}[Вывод по рис.~4]
Знак показателя во втором слагаемом ур.~16 противоречит приложению~B самой
статьи (там фазы входят только через разности $\mathbf r_n-\mathbf r_{n'}$),
независимости от начала координат и ур.~38--39. Рис.~4 и приложение~C
построены по ур.~16 с этим знаком. Наш код использует~\eqref{eq:kd_ad}.
\end{cmpbox}
